\section{A critical-margin upper bound with a remaining iterated logarithm}
\label{sec:critical}
\noindent\textbf{PROVED-ON-PAPER.} The result below is new work in this run,
reviewed by a separate adversarial agent. It is not Lean-checked or externally
reviewed; priority is unresolved. It does not settle the candidate minimax law.
Natural logarithms are used throughout.

Assume the iid hidden-gap Bernoulli experiment of the question. Only requested
coin observations are charged; arithmetic and independent randomness are free.
For known $0<a\le1/2$, assume
\begin{equation}\label{crit:envelope}
 \mu((0,x])\le(x/a)^2\quad(0<x<a).
\end{equation}
The individual gaps, signs, and unknown measurable function $q$ remain hidden.
Only the global condition $\int q\,d\mu\le\rho$ is imposed, with known
$0<\rho\le1/4$. Set $L=1+\log(1/\rho)$ and $H=1+\log L$.

\begin{theorem}\label{crit:theorem}
There is an observable single-coin estimator $Y$, depending only on $a,\rho$,
such that uniformly over \eqref{crit:envelope} and every admissible $q$,
\[
 \mathbb EY=\theta,\qquad
 \mathbb EY^2\le\frac{165}{56}\rho<4\rho,\qquad
 \mathbb ET_Y\le W:=1350a^{-2}LH.
\]
Thus, for every $B>0$, an algorithm of uniform expected work at most $B$ has
worst-case RMS at most
\[
 \min\left\{\rho/2,\sqrt{5400\rho LH/(a^2B)}\right\}.
\]
For $0<\delta\le1/4$, let $k$ be the least odd integer at least
$8\log(1/\delta)$. A confidence radius
\[
 \min\{\rho/2,4\sqrt{\rho Wk/B}\}
\]
is attainable with the same uniform expected-work budget. These statements
give no deterministic work cap or high-probability work bound.
\end{theorem}

\subsection{An entirely observable finite classifier}
For $j\ge0$, put
\[
 \varepsilon_j=16^{-(j+1)},\quad
 N_j=\left\lceil(a^2\rho\varepsilon_j)^{-1}\right\rceil,\quad
 K_j=\left\lceil\log_2(a\sqrt{N_j})\right\rceil.
\]
For $0\le k\le K_j$, let $r_k=a2^{-(k+1)}$. If $H_n$ is the number of
observed heads after $n$ flips of the current coin, define
\[
 L^+_{k,n}=(1+2r_k)^{H_n}(1-2r_k)^{n-H_n},\qquad
 L^-_{k,n}=(1-2r_k)^{H_n}(1+2r_k)^{n-H_n}.
\]
At each $n<N_j$, output zero and stop if some $L^-_{k,n}$ reaches
$(K_j+1)/\varepsilon_j$, or output one and stop if some $L^+_{k,n}$ reaches
$(K_j+1)/(\rho\varepsilon_j)$. Give zero priority in a simultaneous crossing.
If neither decision has occurred before $N_j$, output
$\mathbf1_{\{H_{N_j}>N_j/2\}}$. Call the output $Z_j$ and work $T_j$.
Only finitely many likelihoods are checked at any time.

Under the wrong sign the corresponding likelihood is a nonnegative
supermartingale with expected multiplier $1-4r_kd\le1$. Applying the
bounded-horizon crossing inequality to each of the $K_j+1$ likelihoods gives
wrong-crossing probability at most $\varepsilon_j$ on positives and
$\rho\varepsilon_j$ on negatives. The wrong forced-majority event is a
subset of the full-tape majority-error event, whose conditional probability
is at most $\exp(-2N_jd^2)$ by Hoeffding. No independence between stopping
and this event is used. Stochastic dominance from \eqref{crit:envelope} yields
\[
 \mathbb E_\mu e^{-2ND^2}
 \le\frac{1-e^{-2Na^2}}{2Na^2}\le\frac1{2Na^2}.
\]
Consequently, for every admissible $q$,
\begin{equation}\label{crit:error}
 \mathbb P(Z_j\ne S)\le\tfrac52\rho\varepsilon_j,
 \qquad
 \mathbb P(Z_j=1,S=0)\le\tfrac32\rho\varepsilon_j.
\end{equation}
The $\theta\varepsilon_j$ positive error was bounded by
$\rho\varepsilon_j$; the assumption was never strengthened to $q(d)\le\rho$.
For completeness, the concentration input has a direct Bernoulli proof:
the centered log moment-generating function $\phi(\lambda)$ has
$\phi(0)=\phi'(0)=0$ and $\phi''(\lambda)\le1/4$, because the tilted
Bernoulli variance is at most $1/4$. Hence $\phi(\lambda)\le\lambda^2/8$.
Multiply these bounds for independent observations and apply Markov's
inequality with $\lambda=4d$ (or $-4d$) to obtain $e^{-2Nd^2}$.
This argument also applies to the independent failure indicators used below.

\subsection{Conditional work and global-rarity rearrangement}
Write $x=\min(d,a)$ and
\[
 A_s=1+\log(K_j+1)+\log(1/\varepsilon_j)+s\log(1/\rho),\quad s=0,1.
\]
If $x\ge N_j^{-1/2}$, choose $k=\lceil\log_2(a/x)\rceil\le K_j$.
Then $x/4\le r_k\le x/2$, $r_k\le1/4$, and the correct-sign log-likelihood
has iid increments with mean
\[
 g=\tfrac12\log(1-4r_k^2)
       +d\log\frac{1+2r_k}{1-2r_k}
 \ge-\tfrac83r_k^2+4dr_k\ge\tfrac23x^2.
\]
Here $\log(1-z)\ge-z/(1-z)$ and
$\log((1+u)/(1-u))\ge2u$ suffice. Every upward increment is less than one.
For the crossing time $\sigma$ of the correct threshold, the bounded
stopped-sum identity at $\sigma\wedge n$ gives
$g\mathbb E(\sigma\wedge n)\le A_s$. Monotone convergence therefore gives
$\mathbb E\sigma\le(3/2)A_s/x^2$. The capped classifier work is at most
$\min(N_j,\sigma)$, including when the crossing occurs exactly at the cap.
If $x<N_j^{-1/2}$ the same conclusion below holds trivially because $A_s\ge1$:
\begin{equation}\label{crit:condwork}
 \mathbb E(T_j\mid d,s)\le\min\{N_j,(3/2)A_s/x^2\}.
\end{equation}

Set $c=3/2$. For the negative-sign contribution stochastic dominance gives
\[
 \mathbb E\min\{N,cA_0/x^2\}
 \le cA_0a^{-2}[1+\log_+(Na^2)].
\]
For positives, layer cake and
$\mathbb E[q\mathbf1_E]\le\min\{\rho,\mu(E)\}$ give
\begin{align*}
 \mathbb E\bigl[q(D)\min\{N,cA_1/x^2\}\bigr]
 &\le cA_1a^{-2}[1+\log_+(Na^2\rho/(cA_1))]\\
 &\le cA_1a^{-2}[1+\log_+(Na^2\rho)].
\end{align*}
Indeed the sharp upper is the integral of
$\min\{N,cA_1/(a^2u)\}$ over $0<u<\rho$. If
$Na^2\rho\le cA_1$ this is $N\rho\le cA_1/a^2$; otherwise it is exactly
$(cA_1/a^2)[1+\log(Na^2\rho/(cA_1))]$. This proof also covers envelope laws
with atoms, through layer cake rather than a quantile-density assumption.
Bounding the negative part by the full unweighted integral and adding the
positive part bounds $\mathbb ET_j$.

The following loose universal bounds follow directly from the definitions:
\[
 \begin{gathered}
 N_ja^2\le2/(\rho\varepsilon_j),\qquad K_j+1\le5(L+j+1),\\
 A_0\le7H(j+1),\qquad A_1\le7H(j+1)+L,\\
 1+\log(N_ja^2)\le5L(j+1),\qquad
 1+\log(N_ja^2\rho)\le5(j+1).
 \end{gathered}
\]
For example, use $\log(L+j+1)\le\log L+j+1$ and $4\log2<3$ in the
bound on $A_0$. Since $L,H\ge1$, substitution proves
\begin{equation}\label{crit:work}
 \mathbb ET_j\le(225/2)a^{-2}LH(j+1)^2.
\end{equation}

\subsection{Randomized telescoping and exact expected-budget aggregation}
Independently draw $J$ with $p_j=2^{-(j+1)}$, then draw a fresh coin. For
$J=0$, set $\Delta_0=Z_0$. For $J=j\ge1$, run $Z_j$ and $Z_{j-1}$ on
two successive disjoint observation blocks of that \emph{same} coin and set
$\Delta_j=Z_j-Z_{j-1}$. Return $Y=\Delta_J/p_J$.

Conditional on $(D,S)$, the observations after the first bounded stopping
time again form an iid Bernoulli tape with the original parameter. To verify
this directly, condition on each possible first stopping time $n$ and each
observed length-$n$ prefix. Its event is prefix-measurable and the unused iid
coordinates have their original law. Summing the finitely many prefix events
proves the suffix assertion. Thus the second block has exactly the classifier
marginal used below; its work can be added in expectation. The two outputs
share their latent sign, which is essential to the disagreement bound.

By \eqref{crit:error},
\[
 \mathbb E\Delta_0^2\le\tfrac{35}{32}\rho,\qquad
 \mathbb E\Delta_j^2\le\tfrac{85}{32}\rho16^{-j}\quad(j\ge1).
\]
For the latter, different decisions imply that at least one is wrong about
the shared sign. Hence Tonelli gives
\[
 \mathbb EY^2\le\frac{35}{16}\rho+
      \frac{85}{16}\rho\sum_{j\ge1}8^{-j}
   =\frac{165}{56}\rho<4\rho.
\]
The same bounds give $\sum_j\mathbb E|\Delta_j|<\infty$; signed
expectations may therefore be interchanged and telescoped, giving
$\mathbb EY=\lim_j\mathbb EZ_j=\theta$. This proves integrability before
using the signed telescope. The expected work of level $j$ is at most twice
\eqref{crit:work}, and
$\sum_jp_j(j+1)^2=6$. Thus $\mathbb ET_Y\le1350a^{-2}LH$.
Every call terminates almost surely, and the resulting expected work is finite.

For budget $B$, set $\lambda=B/W$, independently draw
$M\sim\operatorname{Poisson}(\lambda)$, and average the independent fresh-coin
outputs as $\lambda^{-1}\sum_{i=1}^M Y_i$. Conditioning on $M$ proves
expectation $\theta$, variance $\mathbb EY^2/\lambda\le4\rho W/B$, and
expected work at most $B$. This remains valid for $\lambda<1$. Project onto
$[0,\rho]$, which improves squared error without preserving unbiasedness.
If the stated RMS bound exceeds $\rho/2$, choose the zero-work constant
$\rho/2$ instead. This proves the RMS assertion.

For confidence, run $k$ independent copies with budgets $B/k$ and take their
median. Each has failure probability at most $1/4$ at radius
$4\sqrt{\rho Wk/B}$ by Chebyshev. Hoeffding's inequality for the independent
failure indicators bounds median failure by $e^{-k/8}\le\delta$. If
$\rho/2$ is the smaller displayed radius, explicitly choose the constant
estimator instead. Total expected work is at most $B$.

\subsection{Exact critical candidate and the unresolved gap}
For the exact law $F(d)=(d/a)^2$, direct integration yields
\[
 \rho=(t/a)^2[1+2\log(a/t)],\qquad
 \Psi(h)=\begin{cases}
 [2ht(1+\log(a/t))-h^2]/a^2,&0<h<t,\\
 \rho,&h\ge t.
 \end{cases}
\]
Set $u=1+2\log(a/t)$. Then $\rho=u e^{1-u}$ and
$L=u-\log u$, whence $L\le u\le2L$. For $h<t$,
$\Psi(h)/h=(t/a^2)(u+1-h/t)$ lies between $tu/a^2$ and
$t(u+1)/a^2$. Since $tu/a^2=\sqrt{\rho u}/a$, including the transition
at $h=t$ proves
\[
 \Psi(h)\asymp\min\{\rho,h\sqrt{\rho L}/a\}
\]
with universal constants, jointly in $\rho,h,a$. Conditional on the audited
arbitrary-law lower theorem, Theorem~\ref{crit:theorem} therefore leaves a
factor at most $\sqrt{1+\log L}$. It improves the critical upper bracket but
does not remove this factor, prove it necessary, or establish arbitrary-law
optimality. At fixed positive $\rho$, the construction gives $B^{-1/2}$ RMS;
this fixed-rarity statement must not be confused with a uniform rare-event
matching theorem.

\section{A finite, arbitrary-law upper bracket}
\noindent\textbf{PROVED-ON-PAPER; established allocation principle.}
The classifier analysis also gives a constructive upper bracket for every
law on $(0,1/2]$, including atoms, singular laws, and infinite inverse moments.
No comparison of this bracket to $\Psi$ is asserted.

Set $a=1/2$. Choose finitely many integer caps $N_j\ge4$ and tolerances
$0<\varepsilon_j\le1$, $0\le j\le m$, and use the preceding classifiers.
Let
\[
 \beta_j=\mathbb E_\mu e^{-2N_jD^2},\quad
 e_j=2\rho\varepsilon_j+\beta_j,\quad
 A_{j,s}=1+\log(K_j+1)+\log(1/\varepsilon_j)+s\log(1/\rho).
\]
For a bounded nonnegative measurable function $f$, define
\[
 Q_\rho(f)=\int_0^\infty\min\{\rho,\mu(f>v)\}\,dv.
\]
Layer cake gives $\int qf\,d\mu\le Q_\rho(f)$ for every feasible $q$.
Equality follows by choosing the largest $\rho$ mass of $f$, with a fractional
value of $q$ on a cutoff atom. Thus this formula does not require a nonatomic
law. Put
\[
 f_{j,s}(d)=\min\{N_j,(3/2)A_{j,s}/d^2\},\qquad
 c_j=\mathbb E f_{j,0}+Q_\rho(f_{j,1}).
\]
The preceding conditional-work proof gives a uniform bound $c_j$; all these
integrals are finite because the functions are bounded. Let
\[
 v_0=\rho+\rho\varepsilon_0+\beta_0,\quad C_0=c_0,
 \qquad v_j=e_j+e_{j-1},\quad C_j=c_j+c_{j-1}\quad(j\ge1).
\]
Use independent copies of the same-coin increments, independently across
levels. With $S_m=\sum_{j=0}^m\sqrt{v_jC_j}$, independently take
Poisson numbers of copies at means
$\lambda_j=B\sqrt{v_j/C_j}/S_m$. Their sum, with level $j$ divided by
$\lambda_j$, has expectation $\mathbb EZ_m$, variance at most $S_m^2/B$,
and expected work at most $B$. Therefore
\[
 R_{\mu,\rho}(B)\le
 \min\{\rho/2,\sqrt{e_m^2+S_m^2/B}\}.
\]
All constants here are positive. The RMS bound uses the exact bias--variance
identity, and clipping onto $[0,\rho]$ cannot increase risk.

The known law can determine these finitely many design constants, or certified
upper bounds on them, before sampling. This gives no access to the magnitude
of a sampled coin. For an abstract probability measure this is existence at
the fixed known law; uniform computability from every representation of a
measure is not claimed. Floating-point quadrature alone supplies numerical
design evidence, not certified risk constants. The Poisson allocation and
randomized telescoping principles are established methods; no novelty is
claimed for them.

\section{A reusable bounded-channel upper principle}
\noindent\textbf{PROVED-ON-PAPER; no lower transfer.}
The critical theorem and arbitrary-law upper bracket hold with the same
constants if each hidden state produces conditionally iid observations
$X\in[0,1]$ of mean $1/2+(2S-1)D$. Charge one unit per observation; the
algorithm receives no hidden magnitude, sign, or selected-state access.
Replace the reference processes by
\[
 M^\pm_{k,n}=\prod_{i=1}^n[1\pm2r_k(2X_i-1)]
\]
and keep the same thresholds, caps, and empirical-mean decision. The factors
are positive. Under the wrong sign their conditional mean is $1-4r_kd\le1$.
For the correct sign, concavity of the logarithm gives
\[
 \log(1+2r(2X-1))\ge
 X\log(1+2r)+(1-X)\log(1-2r).
\]
Its expected log increment is therefore at least the Bernoulli drift used
above; upward increments remain below one. For the concentration step the
second derivative of a centered log moment-generating function is the variance
under a tilted law on $[0,1]$, at most $1/4$. Hence the previous mgf proof
applies. All conditional work, global-rarity integration, and telescoping
arguments carry over unchanged. This rule directly uses the real observations.

For a concrete derived synthetic example, take, conditional on $d,s$,
\[
 X\sim\operatorname{Unif}
 [1/2+(2s-1)d-1/4,\;1/2+(2s-1)d+1/4],
 \quad 0<d\le a\le1/4.
\]
With independent noises on revisits this is a valid bounded channel with
variance $1/48$. The quadratic-margin upper theorem applies, charging one
real draw per step. No minimax equality or empirical application is asserted.

An exact information-discarding reduction is also available: using an
independent $U\sim\operatorname{Unif}[0,1]$, the variable
$\mathbf1_{\{U\le X\}}$ is a Bernoulli observation with the required
conditional mean. One real observation is charged per converted bit and the
target is unchanged. This transfers Bernoulli upper algorithms; it does not
transfer a Bernoulli lower bound to every channel. If instead $X$ is
deterministically its conditional mean, one observation reveals $S$, and $n$
fresh states yield variance at most $\rho/n$ without an inverse-gap factor.
Nothing here extends a bounded-observation theorem to unbounded observations.

Extra cross-state structure also changes the question. As an explicit
different experiment, suppose the outer state $Z\sim\operatorname{Unif}[-1,1]$
is observed, $\beta\in[-1/2,1/2]$ is unknown, and
\[
 X=\tfrac12+(Z-\beta)/8+\xi,\qquad
 \xi\sim\operatorname{Unif}[-1/16,1/16].
\]
Independent pairs have $X\in[0,1]$, and the positive-mean fraction is
$\theta=(1-\beta)/2$. The average of $Z-8(X-1/2)$ is unbiased for $\beta$
with variance $1/(12n)$; the corresponding estimate of $\theta$ has variance
$1/(48n)$. Charging one unit per outer draw and one per inner observation
gives work $2n$. The margin near zero is linear, yet this common regression
parameter permits $B^{-1/2}$ RMS without a margin logarithm. Observed states
and cross-state structure are explicit extra information, so this calculation
does not contradict the hidden unstructured reservoir lower theorem.
